\chapter{2000年上海卷（理）}

\section{填空题}

\begin{problem}
已知向量 \(\overrightarrow{OA} = (-1, 2)\)，\(\overrightarrow{OB} = (3, m)\)，若 \(\overrightarrow{OA} \perp \overrightarrow{OB}\)，则 \(m = \)\fillinblank{}.
\end{problem}

\begin{answer}
\(\dfrac{3}{2}\)
\end{answer}

\begin{solution}
由垂直向量的数量积为零，得 \((-1,2)\cdot(3,m)=-3+2m=0\)，所以 \(m=\dfrac{3}{2}\).
\end{solution}

\begin{problem}
函数 \( y = \log_2 \frac{2x - 1}{3 - x} \) 的定义域为\fillinblank{}.
\end{problem}

\begin{answer}
\(\left(\frac{1}{2},3\right)\)
\end{answer}

\begin{solution}
对数的真数必须为正，故 \(\dfrac{2x-1}{3-x}>0\). 临界点为 \(x=\dfrac{1}{2}\) 和 \(x=3\)：当 \(x<\dfrac{1}{2}\) 时分子为负、分母为正，真数为负；当 \(\dfrac{1}{2}<x<3\) 时分子、分母均为正，真数为正；当 \(x>3\) 时分子为正、分母为负，真数为负. 因此定义域为 \(\left(\dfrac{1}{2},3\right)\).
\end{solution}

\begin{problem}
圆锥曲线 \(\begin{cases}
x = 4 \sec \theta + 1,\\
y = 3 \tan \theta,
\end{cases}\) 的焦点坐标是\fillinblank{}.
\end{problem}

\begin{answer}
\((-4,0)\)，\((6,0)\)
\end{answer}

\begin{solution}
由 \(x-1=4\sec\theta\)、\(y=3\tan\theta\)，消去参数得 \(\dfrac{(x-1)^2}{16}-\dfrac{y^2}{9}=1\). 这是中心为 \((1,0)\) 的双曲线，其中 \(a=4\)、\(b=3\)，故 \(c=\sqrt{a^2+b^2}=5\). 两焦点为 \((1-5,0)=(-4,0)\) 和 \((1+5,0)=(6,0)\).
\end{solution}

\begin{problem}
计算： \( \displaystyle\lim_{n\to\infty}\left(\frac{n}{n+2}\right)^n = \)\fillinblank{}.
\end{problem}

\begin{answer}
\(\e^{-2}\)
\end{answer}

\begin{solution}
\(\left(\dfrac{n}{n+2}\right)^n=\left(1+\dfrac{2}{n}\right)^{-n}\). 由 \(\displaystyle\lim_{n\to\infty}\left(1+\dfrac{2}{n}\right)^n=\e^2\)，可得原极限为 \(\e^{-2}\).
\end{solution}

\begin{problem}
已知 \( f(x) = 2^{x} + b \) 的反函数为 \( f^{-1}(x) \)，若 \( y = f^{-1}(x) \) 的图象经过点 \( Q(5,2) \)，则 \(b=\)\fillinblank{}.
\end{problem}

\begin{answer}
\(1\)
\end{answer}

\begin{solution}
由反函数图象关于直线 \(y=x\) 对称，点 \(Q(5,2)\) 在 \(y=f^{-1}(x)\) 上等价于点 \((2,5)\) 在 \(y=f(x)\) 上. 因而 \(5=2^2+b=4+b\)，所以 \(b=1\).
\end{solution}

\begin{problem}
根据上海市人大十一届三次会议上的市政府工作报告，1999 年上海市完成 \(GDP\)（\(GDP\) 是指国内生产总值）\(4035\text{ 亿元}\)，2000 年上海市 \(GDP\) 预期增长 \(9\%\). 市委、市政府提出本市常住人口每年的自然增长率将控制在 \(0.08\%\)，若 \(GDP\) 与人口均按这样的速度增长，则要使本市年人均 \(GDP\) 达到或超过 1999 年的 \(2\text{ 倍}\)，至少需要\fillinblank{}\(\text{年}\). 按：1999 年本市常住人口总数约 \(1300\text{ 万}\).
\end{problem}

\begin{answer}
\(9\)
\end{answer}

\begin{solution}
设经过 \(n\) 年后，\(GDP\) 与人口分别为初始值的 \(1.09^n\) 和 \(1.0008^n\) 倍，则人均 \(GDP\) 为初始值的 \(\left(\dfrac{1.09}{1.0008}\right)^n\) 倍. 因此需要
\(\left(\dfrac{1.09}{1.0008}\right)^n\geqslant2\). 由于 \(\dfrac{1.09}{1.0008}\approx1.08913>1\)，左端随 \(n\) 增大而增大；直接计算得 \(\left(\dfrac{1.09}{1.0008}\right)^8\approx1.9799<2\)，而 \(\left(\dfrac{1.09}{1.0008}\right)^9\approx2.1563>2\)，所以至少需要 \(9\) 年.
\end{solution}

\begin{problem}
命题 \(A\)：底面为正三角形，且顶点在底面的射影为底面中心的三棱锥是正三棱锥. 命题 \(A\) 的等价命题 \(B\) 可以是：底面为正三角形，且\fillinblank{}的三棱锥是正三棱锥.
\end{problem}

\begin{answer}
侧棱相等
\end{answer}

\begin{solution}
设底面正三角形的中心为 \(O\)，顶点为 \(P\). 若 \(PO\perp\) 底面，则 \(OA=OB=OC\)，由勾股定理得 \(PA=PB=PC\). 反之，若 \(PA=PB=PC\)，则点 \(P\) 在底面各边的垂直平分面上；这些垂直平分面的交线是过 \(O\) 且垂直于底面的直线，所以 \(P\) 在该直线上，即顶点在底面的射影为 \(O\). 因而等价条件可以是侧棱相等.
\end{solution}

\begin{problem}
设函数 \( y = f(x) \) 是最小正周期为 2 的偶函数，它在区间 \([0,1]\) 上的图象为如图所示的线段 \(AB\)，则在区间 \([1,2]\) 上 \( f(x) = \)\fillinblank{}.

\bitmapfigure[width=4.0cm]{img/2000/shanghai_science/q8_fig1.png}
\end{problem}

\begin{answer}
\(x\)
\end{answer}

\begin{solution}
在 \([0,1]\) 上，线段 \(AB\) 经过 \(A(0,2)\)、\(B(1,1)\)，故 \(f(t)=2-t\)（\(0\leqslant t\leqslant1\)）. 对 \(1\leqslant x\leqslant2\)，由偶性和周期为 \(2\) 得 \(f(x)=f(-x)=f(2-x)=2-(2-x)=x\).
\end{solution}

\begin{problem}
在二项式 \((x - 1)^{11}\) 的展开式中，系数最小的项的系数为\fillinblank{}. （结果用数值表示）
\end{problem}

\begin{answer}
\(-462\)
\end{answer}

\begin{solution}
展开式通项为 \(T_{r+1}=\mathrm C_{11}^{r}x^{11-r}(-1)^r\)（\(0\leqslant r\leqslant11\)）. 负系数对应奇数 \(r\)，其中 \(\mathrm C_{11}^{1}=11\)、\(\mathrm C_{11}^{3}=165\)、\(\mathrm C_{11}^{5}=462\)、\(\mathrm C_{11}^{7}=330\)、\(\mathrm C_{11}^{9}=55\)、\(\mathrm C_{11}^{11}=1\)，最大负系数的绝对值为 \(462\)，所以最小系数为 \(-462\).
\end{solution}

\begin{problem}
红、黄、蓝三种颜色的旗帜各3面，在每种颜色的3面旗帜上分别标上号码1、2和3，现任取出3面，它们的颜色与号码均不相同的概率是\fillinblank{}.
\end{problem}

\begin{answer}
\(\frac{1}{14}\)
\end{answer}

\begin{solution}
从 \(9\) 面旗帜中任取 \(3\) 面，共有 \(\mathrm C_{9}^{3}=84\) 种等可能取法. 要使颜色和号码均不相同，先选定三种颜色，再将号码 \(1\)，\(2\)，\(3\) 分配给红、黄、蓝三色，合适取法有 \(3!=6\) 种. 所求概率为 \(\dfrac{6}{84}=\dfrac{1}{14}\).
\end{solution}

\begin{problem}
在极坐标系中，若过点 \((3,0)\) 且与极轴垂直的直线交曲线 \(\rho = 4\cos \theta\) 于 \(A\)，\(B\) 两点，则 \(|AB| =\)\fillinblank{}.
\end{problem}

\begin{answer}
\(2\sqrt{3}\)
\end{answer}

\begin{solution}
曲线 \(\rho=4\cos\theta\) 化为直角坐标方程 \(x^2+y^2=4x\)，即 \((x-2)^2+y^2=4\). 过 \((3,0)\) 且垂直于极轴的直线为 \(x=3\)，与圆联立得 \(1+y^2=4\)，所以交点纵坐标为 \(y=\pm\sqrt{3}\). 因而 \(|AB|=2\sqrt{3}\).
\end{solution}

\begin{problem}
在等差数列 \(\{a_{n}\}\) 中，若 \(a_{10} = 0\)，则有等式 \(a_{1} + a_{2} + \cdots + a_{n} = a_{1} + a_{2} + \cdots + a_{19-n}\)（\(n < 19\)，\(n \in \N^*\)）成立，类比上述性质，相应地：在等比数列 \(\{b_{n}\}\) 中，若 \(b_{9} = 1\)，则有等式\fillinblank{}成立.
\end{problem}

\begin{answer}
\(b_{1}b_{2}\cdots b_{n}=b_{1}b_{2}\cdots b_{17-n}\)，\(n<17\)，\(n\in\N^*\)
\end{answer}

\begin{solution}
设等比数列首项为 \(b_1\)，公比为 \(q\). 由 \(b_9=b_1q^8=1\) 可知 \(q\neq0\)，从而
\(b_1=q^{-8}\)，且前 \(n\) 项积为 \(b_1b_2\cdots b_n=b_1^nq^{1+2+\cdots+(n-1)}=q^{-8n+\frac{n(n-1)}{2}}=q^{\frac{n(n-17)}{2}}\). 同理，\(b_1b_2\cdots b_{17-n}=q^{\frac{(17-n)((17-n)-17)}{2}}=q^{\frac{n(n-17)}{2}}\). 因此对所有 \(n<17\)、\(n\in\N^*\)，都有 \(b_1b_2\cdots b_n=b_1b_2\cdots b_{17-n}\).
\end{solution}

\section{单选题}

\begin{problem}
复数 \(z = -3\left(\cos \frac{\pi}{5} - \i\sin \frac{\pi}{5}\right)\)（\(\i\) 是虚数单位）的三角形式是（\quad）

\choices
    {\(3\left[\cos \left(-\frac{\pi}{5}\right) + \i\sin \left(-\frac{\pi}{5}\right)\right]\)}
    {\(3\left(\cos\frac{\pi}{5}+\i\sin\frac{\pi}{5}\right)\)}
    {\(3\left(\cos\frac{4\pi}{5}+\i\sin\frac{4\pi}{5}\right)\)}
    {\(3\left(\cos \frac{6\pi}{5} -\i\sin \frac{6\pi}{5}\right)\)}
\end{problem}

\begin{answer}
C
\end{answer}

\begin{solution}
因为 \(-3\cos\dfrac{\pi}{5}=3\cos\dfrac{4\pi}{5}\)，且 \(3\sin\dfrac{\pi}{5}=3\sin\dfrac{4\pi}{5}\)，所以
\(z=3\left(\cos\dfrac{4\pi}{5}+\i\sin\dfrac{4\pi}{5}\right)\). 选项 D 代数上也表示同一个复数，但未写成标准的 \(r\left(\cos\theta+\i\sin\theta\right)\) 形式，故按题目要求选 C.
\end{solution}

\begin{problem}
设有不同直线 \(a\)，\(b\) 和不同的平面 \(\alpha\)，\(\beta\)，\(\gamma\)，给出下列三个命题：
\begin{circlelist}
\item 若 \(a \parallel \alpha\)，\(b \parallel \alpha\)，则 \(a \parallel b\)；
\item 若 \(a \parallel \alpha\)，\(a \parallel \beta\)，则 \(\alpha \parallel \beta\)；
\item 若 \(\alpha \perp \gamma\)，\(\beta \perp \gamma\)，则 \(\alpha \parallel \beta\).
\end{circlelist}
其中正确的个数是（\quad）

\choices
    {\(0\)}
    {\(1\)}
    {\(2\)}
    {\(3\)}
\end{problem}

\begin{answer}
A
\end{answer}

\begin{solution}
第一个命题不正确：两条都平行于同一平面的直线不一定互相平行，例如它们可以相交. 第二个命题不正确：同一直线可以同时平行于两个相交平面，例如取两个相交的坐标平面，并取一条与它们的交线平行且不在任一平面内的直线. 第三个命题也不正确：垂直于同一平面的两个平面不一定平行，例如两个不同的竖直平面都垂直于水平面. 因而正确命题的个数为 \(0\)，选 A.
\end{solution}

\begin{problem}
若集合 \( S = \{y \mid y = 3^x, x \in \R\} \)，\( T = \{y \mid y = x^2 - 1, x \in \R\} \)，则 \( S \cap T \) 是（\quad）

\choices
    {\(S\)}
    {\(T\)}
    {\(\varnothing\)}
    {有限集}
\end{problem}

\begin{answer}
A
\end{answer}

\begin{solution}
因为 \(3^x>0\) 且当 \(x\) 取遍实数时 \(3^x\) 取遍所有正数，所以 \(S=(0,+\infty)\). 又 \(x^2\geqslant0\)，故 \(T=[-1,+\infty)\). 因而 \(S\subseteq T\)，从而 \(S\cap T=S\)，选 A.
\end{solution}

\begin{problem}
下列命题中正确的命题是（\quad）

\choices
    {若点 \(P(a, 2a)\)（\(a \neq 0\)）为角 \(\alpha\) 终边上一点，则 \(\sin \alpha = \frac{2\sqrt{5}}{5}\)}
    {同时满足 \(\sin \alpha = \frac{1}{2}\)，\(\cos \alpha = \frac{\sqrt{3}}{2}\) 的角 \(\alpha\) 有且只有一个}
    {当 \(|a| < 1\) 时，\(\tan (\arcsin a)\) 的值恒正}
    {三角方程 \(\tan \left(x + \frac{\pi}{3}\right) = \sqrt{3}\) 的解集为 \(\{x \mid x = k\pi, k \in \Z\}\)}
\end{problem}

\begin{answer}
D
\end{answer}

\begin{solution}
若 \(P(a,2a)\) 在角 \(\alpha\) 的终边上，则 \(r=\sqrt{a^2+4a^2}=|a|\sqrt{5}\)，所以 \(\sin\alpha=\dfrac{2a}{|a|\sqrt{5}}\)，当 \(a<0\) 时为负，命题 \(A\) 不正确. 命题 \(B\) 中的正弦、余弦只确定角的同余类，所有 \(\alpha=\dfrac{\pi}{6}+2k\pi\)（\(k\in\Z\)）都满足条件，并非只有一个角. 当 \(|a|<1\) 时，\(\arcsin a\in\left[-\dfrac{\pi}{2},\dfrac{\pi}{2}\right]\)，且 \(\tan(\arcsin a)=\dfrac{a}{\sqrt{1-a^2}}\)，其符号随 \(a\) 改变，命题 C 不正确. 对命题 D，\(\tan u=\sqrt{3}\) 的通解为 \(u=\dfrac{\pi}{3}+k\pi\)，令 \(u=x+\dfrac{\pi}{3}\)，得 \(x=k\pi\)（\(k\in\Z\)），故 D 正确.
\end{solution}

\section{解答题}

\begin{problem}
已知椭圆 \(C\) 的焦点分别为 \(F_{1}(-2\sqrt{2},0)\) 和 \(F_{2}(2\sqrt{2},0)\). 长轴长为 6，设直线 \(y = x + 2\) 交椭圆 \(C\) 于 \(A\)，\(B\) 两点，求线段 \(AB\) 的中点坐标.
\end{problem}

\begin{solution}
设椭圆的标准方程为 \(\dfrac{x^2}{a^2}+\dfrac{y^2}{b^2}=1\). 由长轴长为 \(6\) 得 \(a=3\)，由焦点坐标得 \(c=2\sqrt{2}\)，从而 \(b^2=a^2-c^2=9-8=1\). 所以椭圆方程为 \(\dfrac{x^2}{9}+y^2=1\).

将直线 \(y=x+2\) 代入椭圆方程，得 \(\dfrac{x^2}{9}+(x+2)^2=1\)，即 \(10x^2+36x+27=0\). 其判别式为 \(36^2-4\cdot10\cdot27=216>0\)，故确有两个交点. 设两交点横坐标为 \(x_1\)，\(x_2\)，由韦达定理，\(x_1+x_2=-\dfrac{36}{10}=-\dfrac{18}{5}\). 交点纵坐标分别为 \(x_1+2\)，\(x_2+2\)，所以 \(AB\) 的中点为
\(\left(\dfrac{x_1+x_2}{2},\dfrac{(x_1+2)+(x_2+2)}{2}\right)=\left(-\dfrac{9}{5},\dfrac{1}{5}\right)\).
\end{solution}

\begin{problem}
如图所示四面体 \(ABCD\) 中，\(AB\)，\(BC\)，\(BD\) 两两互相垂直，且 \(AB = BC = 2\)，\(E\) 是 \(AC\) 中点，异面直线 \(AD\) 与 \(BE\) 所成的角的大小为 \(\arccos \frac{\sqrt{10}}{10}\)，求四面体 \(ABCD\) 的体积.

\bitmapfigure[width=3.7cm]{img/2000/shanghai_science/q18_fig1.png}
\end{problem}

\begin{solution}
以 \(B\) 为原点，分别以 \(\overrightarrow{BA}\)、\(\overrightarrow{BC}\)、\(\overrightarrow{BD}\) 为三个坐标轴正方向，设 \(BD=d\). 则可取 \(A(2,0,0)\)、\(C(0,2,0)\)、\(D(0,0,d)\)，且 \(E\) 是 \(AC\) 的中点，所以 \(E(1,1,0)\). 于是可取异面直线 \(AD\)、\(BE\) 的方向向量为 \(\overrightarrow{AD}=(-2,0,d)\)、\(\overrightarrow{BE}=(1,1,0)\).

设两直线所成的锐角为 \(\varphi\). 由题意 \(\cos\varphi=\dfrac{\sqrt{10}}{10}\)，故
\(\dfrac{|\overrightarrow{AD}\cdot\overrightarrow{BE}|}{|\overrightarrow{AD}|\,|\overrightarrow{BE}|}=\dfrac{2}{\sqrt{d^2+4}\sqrt{2}}=\dfrac{\sqrt{10}}{10}\).
两边平方得 \(\dfrac{2}{d^2+4}=\dfrac{1}{10}\)，所以 \(d^2=16\)，即 \(BD=4\). 由于 \(AB\)、\(BC\)、\(BD\) 两两垂直，四面体体积为
\(V_{ABCD}=\dfrac{1}{6}AB\cdot BC\cdot BD=\dfrac{1}{6}\cdot2\cdot2\cdot4=\dfrac{8}{3}\).
\end{solution}

\begin{problem}
已知函数 \(f(x) = \frac{x^2 + 2x + a}{x}\)，\(x \in [1, +\infty)\).
\begin{enumerate}
\item 当 \(a = \frac{1}{2}\) 时，求函数 \(f(x)\) 的最小值；
\item 若对任意 \(x \in [1, +\infty)\)，\(f(x) > 0\) 恒成立，试求实数 \(a\) 的取值范围.
\end{enumerate}
\end{problem}

\begin{solution}
\begin{enumerate}
\item 当 \(a=\dfrac{1}{2}\) 时，\(f(x)=x+2+\dfrac{1}{2x}\). 对 \(x\geqslant1\)，有 \(f'(x)=1-\dfrac{1}{2x^2}\geqslant\dfrac{1}{2}>0\)，所以 \(f(x)\) 在 \([1,+\infty)\) 上递增，最小值在 \(x=1\) 处取得，为 \(f(1)=1+2+\dfrac{1}{2}=\dfrac{7}{2}\).
\item 因为 \(x\geqslant1>0\)，所以 \(f(x)>0\) 等价于 \(g(x)=x^2+2x+a>0\). 又 \(g'(x)=2x+2>0\)（\(x\geqslant1\)），故 \(g(x)\) 在该区间上递增，最小值为 \(g(1)=a+3\). 因而对所有 \(x\geqslant1\) 恒有 \(f(x)>0\) 当且仅当 \(a+3>0\)，即 \(a>-3\).
\end{enumerate}
\end{solution}

\begin{problem}
根据指令 \((r,\theta)\)（\(r \geqslant 0\)，\(-180^{\circ} \leqslant \theta \leqslant 180^{\circ}\)），机器人在平面上能完成下列动作：先原地旋转角度 \(\theta\)（\(\theta\) 为正时，按逆时针方向旋转 \(\theta\)，\(\theta\) 为负时，按顺时针方向旋转 \(-\theta\)），再朝其面对的方向沿直线行走距离 \(r\).
\begin{enumerate}
\item 现机器人在直角坐标系的坐标原点，且面对 \(x\) 轴正方向，试给机器人下一个指令，使其移动到点 \((4,4)\)；
\item 机器人在完成该指令后，发现在点 \((17,0)\) 处有一小球正向坐标原点作匀速直线滚动，已知小球滚动的速度为机器人直线行走速度的 2 倍，若忽略机器人原地旋转所需的时间，问机器人最快可在何处截住小球？并给出机器人截住小球所需的指令.（结果精确到小数点后两位）
\end{enumerate}

\bitmapfigure[width=4.6cm]{img/2000/shanghai_science/q20_fig1.png}
\end{problem}

\begin{solution}
\begin{enumerate}
\item 从原点到 \((4,4)\) 的距离为 \(\sqrt{4^2+4^2}=4\sqrt{2}\)，方向与 \(x\) 轴正方向成 \(45^\circ\)，所以可下指令 \(\left(4\sqrt{2},45^\circ\right)\).
\item 设小球在 \(P(x,0)\) 处被截住，其中 \(0\leqslant x\leqslant17\)，机器人直线行走速度为 \(v\). 小球从 \((17,0)\) 到 \(P\) 所需时间为 \(\dfrac{17-x}{2v}\)，机器人从 \((4,4)\) 到 \(P\) 所需时间为 \(\dfrac{\sqrt{(x-4)^2+16}}{v}\). 两者同时到达时满足
\(\dfrac{\sqrt{(x-4)^2+16}}{v}=\dfrac{17-x}{2v}\).
由于机器人必须不晚于小球到达，需有 \(\dfrac{\sqrt{(x-4)^2+16}}{v}\leqslant\dfrac{17-x}{2v}\). 右端非负，平方后得 \(3x^2+2x-161\leqslant0\)，其根为 \(-\dfrac{23}{3}\) 和 \(7\)，故在 \([0,17]\) 内可行的截球点满足 \(0\leqslant x\leqslant7\). 当 \(x<7\) 时，机器人提前到达后需等待小球，截球时间为 \(\dfrac{17-x}{2v}\)，它随 \(x\) 增大而减小；因此最快截球只能在边界 \(x=7\) 处发生. 此时两者同时到达，机器人行走距离为 \(\sqrt{3^2+(-4)^2}=5\)，截球点为 \(P(7,0)\).

机器人到达 \((4,4)\) 后面对方向为 \(45^\circ\)，而从 \((4,4)\) 指向 \(P(7,0)\) 的方向角为 \(-\arctan\dfrac{4}{3}\approx-53.13^\circ\)，故需要顺时针旋转 \(-53.13^\circ-45^\circ=-98.13^\circ\). 所需指令为 \(\left(5,-98.13^\circ\right)\).
\end{enumerate}
\end{solution}

\begin{problem}
在 \(xOy\) 平面上有一点列 \(P_{1}(a_{1},b_{1})\)，\(P_{2}(a_{2},b_{2})\)，\(\cdots\)，\(P_{n}(a_{n},b_{n})\)，\(\cdots\)，对每个自然数 \(n\)，点 \(P_{n}\) 位于函数 \(y = 2000 \cdot \left(\frac{a}{10}\right)^{n}\)（\(0 < a < 10\)）的图象上，且点 \(P_{n}\)，点 \((n,0)\) 与点 \((n + 1,0)\) 构成一个以 \(P_{n}\) 为顶点的等腰三角形.
\begin{enumerate}
\item 求点 \(P_{n}\) 的纵坐标 \(b_{n}\) 的表达式；
\item 若对每个自然数 \(n\)，以 \(b_{n}\)，\(b_{n+1}\)，\(b_{n+2}\) 为边长能构成一个三角形，求 \(a\) 取值范围；
\item 设 \(B_{n} = b_{1}b_{2}\cdots b_{n}\)（\(n \in \N\)），若 \(a\) 取（2）中确定的范围内的最小整数，求数列 \(\{B_{n}\}\) 的最大项的项数.
\end{enumerate}
\end{problem}

\begin{solution}
\begin{enumerate}
\item 因为 \(P_n\) 与 \((n,0)\)、\((n+1,0)\) 构成以 \(P_n\) 为顶点的等腰三角形，所以 \(P_n\) 到底边两个端点的距离相等. 设 \(P_n=(a_n,b_n)\)，则
\((a_n-n)^2+b_n^2=(a_n-n-1)^2+b_n^2\)，化简得 \(2a_n-2n-1=0\)，所以 \(a_n=n+\dfrac{1}{2}\). 又 \(P_n\) 在函数 \(y=2000\left(\dfrac{a}{10}\right)^n\) 的图象上，故
\(b_n=2000\left(\dfrac{a}{10}\right)^n\)，其中 \(b_n>0\).
\item 记 \(r=\dfrac{a}{10}\in(0,1)\). 由 \(b_n=2000r^n\)，三项 \(b_n,b_{n+1},b_{n+2}\) 能构成三角形当且仅当 \(b_n<b_{n+1}+b_{n+2}\)，即 \(1<r+r^2\). 因此 \(r>\dfrac{\sqrt5-1}{2}\)，从而 \(a\in\left(5(\sqrt5-1),10\right)\).
\item 区间 \(\left(5(\sqrt5-1),10\right)\) 内的最小整数为 \(a=7\)，记 \(r=0.7\). 则 \(B_n=\prod_{j=1}^{n}2000r^j=2000^n r^{n(n+1)/2}\)，且 \(\dfrac{B_{n+1}}{B_n}=2000r^{n+1}\). 直接比较相邻项可得 \(B_n\) 在 \(n=21\) 时达到最大，故最大项的项数为 \(21\).
\end{enumerate}
\end{solution}

\begin{problem}
已知复数 \( z_0 = 1 - m\i \)（\( m > 0\)），\( z = x + y\i \) 和 \( \omega = x' + y'\i \)，其中 \(x\)，\(y\)，\(x'\)，\(y'\) 均为实数，\(\i\) 为虚数单位，且对于任意复数 \( z \)，有 \(\omega = \overline{z_0}\cdot\overline{z}\)，\(|\omega| = 2|z|\).
\begin{enumerate}
\item 试求 \(m\) 的值，并分别写出 \(x'\) 和 \(y'\) 用 \(x\)，\(y\) 表示的关系式；
\item 将 \((x, y)\) 作为点 \(P\) 的坐标，\((x', y')\) 作为点 \(Q\) 的坐标，上述关系可以看作是坐标平面上点的一个变换：它将平面上的点 \(P\) 变到这一平面上的点 \(Q\)，当点 \(P\) 在直线 \(y = x + 1\) 上移动时，试求点 \(P\) 经该变换后得到的点 \(Q\) 的轨迹方程；
\item 是否存在这样的直线：它上面的任一点经过上述变换后得到的点仍在该直线上？若存在，试求出所有这些直线；若不存在，则说明理由.
\end{enumerate}
\end{problem}

\begin{solution}
\begin{enumerate}
\item \(\overline{z_0}=1+m\i\)、\(\overline z=x-y\i\)，所以
\(\omega=(1+m\i)(x-y\i)=(x+my)+(mx-y)\i\). 因而 \(x'=x+my\)、\(y'=mx-y\). 又 \(|\omega|=|\overline{z_0}|\,|\overline z|=\sqrt{1+m^2}\,|z|=2|z|\)，故 \(1+m^2=4\). 由 \(m>0\) 得 \(m=\sqrt{3}\)，从而 \(x'=x+\sqrt{3}y\)、\(y'=\sqrt{3}x-y\).
\item 设点 \(P\) 的坐标为 \((t,t+1)\). 由上式，点 \(Q\) 的坐标 \((X,Y)\) 满足 \(X=(1+\sqrt{3})t+\sqrt{3}\)、\(Y=(\sqrt{3}-1)t-1\). 消去 \(t\)，得到
\(Y=(2-\sqrt{3})X+2-2\sqrt{3}\). 因此轨迹方程为 \(y=(2-\sqrt{3})x-2\sqrt{3}+2\).
\item 设所求直线不是竖直直线，方程写为 \(y=kx+b\). 对其上任意点，变换后有 \(x'= (1+\sqrt{3}k)x+\sqrt{3}b\)、\(y'=(\sqrt{3}-k)x-b\). 要使变换后的点仍在原直线上，对一切 \(x\) 必须满足 \(y'=kx'+b\)，比较 \(x\) 的系数和常数项，得
\(\sqrt{3}-k=k(1+\sqrt{3}k)\)，\(-b=\sqrt{3}kb+b\).
前式化为 \(\sqrt{3}k^2+2k-\sqrt{3}=0\)，所以 \(k=\dfrac{\sqrt{3}}{3}\) 或 \(k=-\sqrt{3}\). 对这两个值，后式均给出 \(b=0\). 竖直直线不可能满足条件，因为其上点的 \(x'=x+\sqrt{3}y\) 随 \(y\) 改变. 所有满足条件的直线为 \(y=\dfrac{\sqrt{3}}{3}x\) 和 \(y=-\sqrt{3}x\).
\end{enumerate}
\end{solution}
